\documentclass[10pt]{extarticle}

% ---------- Packages ----------
\usepackage{fancyhdr}
\usepackage{geometry}
\usepackage{hyperref}
\usepackage{xcolor}
\usepackage{amsmath, amsfonts, bm}
\usepackage{subcaption}
\usepackage{graphicx}
\usepackage{tikz}
\usepackage{nomencl}
\usepackage{array}
\usepackage{pdflscape}
\usepackage[export]{adjustbox}
\usepackage[utf8]{inputenc}
\usepackage{parskip}

\usetikzlibrary{shapes.geometric}
\usetikzlibrary{arrows.meta, arrows}

% ---------- Page setup ----------
\geometry{
    a4paper,
    total={170mm,257mm},
    left=20mm,
    top=20mm
}

% ---------- Headers and footers ----------
\pagestyle{fancy}
\pagenumbering{arabic}

\fancyhead[L]{Rocket Physics by Yaroslav Knyazkov}
\fancyfoot[C]{}
\fancyfoot[R]{\thepage}

% ---------- Hyperlinks ----------
\hypersetup{
    colorlinks=true,
    linkcolor={blue!50!black},
    citecolor={blue!50!black},
    urlcolor={blue!80!black}
}

% ---------- Equation numbering ----------
\numberwithin{equation}{subsection}

% ---------- Custom commands ----------
\newcommand{\todo}[1]{%
    \textcolor{red}{(#1)}%
}

\begin{document}
\section{3DOF - Definitions}

\subsection{Assumptions}
\begin{itemize}
    \item rigid body
\end{itemize}

\subsection{Frames}
\begin{itemize}
    \item $\mathcal{F}_0$ - ground
    \item $\mathcal{F}_g$ - gravitational field
    \item $\mathcal{F}_r$ - rocket
    \item $\mathcal{F}_t$ - thrust 
\end{itemize}


\subsection{Forces and moments}
\begin{equation}
    \vec{W} = mg\hat{z}_g, 
\end{equation}
\begin{equation}
    \vec{T} = T\hat{x}_t,
\end{equation}
\begin{equation}
    \vec{A}_f = -A_x\hat{x}_r - A_z\hat{z}_r,
\end{equation}
\begin{equation}
    \vec{A}_m = -A_m\hat{y}_r,
\end{equation}

\subsection{Geometry}
\begin{equation}
    \vec{AB} = -l_B\hat{x}_r - u_B\hat{z}_r.
\end{equation}
\begin{equation}
    \vec{AC} = -l_C\hat{x}_r - u_C\hat{z}_r.
\end{equation}
\begin{equation}
    \vec{AD} = -l_D\hat{x}_r.
\end{equation}

\subsection{Equations of motion}
\begin{gather}
    m\vec{\ddot{x}} = \sum \vec{F} = \vec{W} + \vec{T} + \vec{A}_f \\
    I\ddot{\theta} = \sum \vec{M}(B) \cdot \hat{y}_r = \left[\vec{BC}\times\vec{T} + \vec{BD}\times\vec{A}_f + \vec{A_m}\right]\cdot \hat{y}_r
\end{gather}

\newpage
\section{3DOF - Models}
\subsection{Gravity}
\subsubsection{Magnitude}
\begin{itemize}
    \item constant
    \begin{equation}
        g = g_0
    \end{equation}
    \item radial
    \begin{equation}
        g(r) = g_0 {\left(\frac{R_e}{R_e+r}\right)}^2
    \end{equation}
    \item real data (analytical)
\end{itemize}

\subsubsection{Direction}
\begin{itemize}
    \item parallel
    \begin{equation}
        \hat{z}_g = \hat{z}_0
    \end{equation}
    \item radial (origin at center of earth)
    \begin{equation}
        \hat{z}_g = -\frac{1}{r}\begin{bmatrix}
            x \\ 0 \\ z
        \end{bmatrix}
    \end{equation}
    \item elliptical (origin at center of earth)
\end{itemize}

\subsubsection{Vector field (magnitude + direction)}
data from some dataset

\subsection{Thrust}
\subsubsection{Magnitude}
\begin{equation}
    T = \dot{m}v_e + A_e(p_e - p_{atm}),
\end{equation}
\begin{equation}
    T_{ref}(t) = \mathrm{lerp}\!\left(t; \{t_i\}, \{T_i\}\right)
\end{equation}
\begin{equation}
    T(t, r) = T_{ref}(t) + A_e\left(p_{ref} - p_{amb}(r)\right), \qquad \text{only while } T_{ref}(t) > 0
\end{equation}

\subsubsection{Direction}
done by missaligned or thrust control

\subsection{Mass, center of gravity and moment of inertia}
\subsubsection{Mass}
\begin{equation}
    \dot{m}(t) = -\frac{T_{ref}(t)}{I_{sp}\,g_0}
\end{equation}

\subsubsection{Center of gravity}
\begin{equation}
    l_B(t) = \frac{m_d\,l_d + m_p(t)\,l_p(t)}{m_d + m_p(t)}
\end{equation}

\subsubsection{Moment of inertia}
\begin{equation}
    I(t) = I_d + m_d\bigl(l_B(t)-l_d\bigr)^2 + I_p(t) + m_p(t)\bigl(l_B(t)-l_p\bigr)^2
\end{equation}

\todo{change of mass or moment of inertia due to mechanisms}

\subsection{Aerodynamics}
make them bic C not c
\subsubsection{Magnitudes}
\begin{align}
    A_x &= qSc_x \\
    A_z &= qSc_z \\
    A_m &= qSlc_m
\end{align}

\subsubsection{Dynamic pressure}
mention density, mainly to derive relative air velocity

\subsubsection{Coefficients}
\begin{equation}
    c = c(\alpha, M, Re)
\end{equation}
intoduce mach and reynolds numbers

\todo{Also depends on geometry; control surface extension airbrake, aileron, etc.}
\todo{might want to include pitch rate and for 3d other stuff}
\subsubsection{Linearization of coefficients}
\begin{equation}
    c_z = c_{z,0} + \alpha c_{z,\alpha}
\end{equation}
\begin{equation}
    c_m = c_{m,0} + \alpha c_{m,\alpha}
\end{equation}
\todo{what about $c_d$?}
\begin{equation}
    c_x = c_{x,0} + \frac{c_z^2}{e\pi\lambda}
\end{equation}
(thats for airfoils, what is in general?)

\subsection{Atmosphere}
\subsubsection{Models}
present what data is given

\subsubsection{Pressure}
\subsubsection{Density}
+ humid density derivation

\subsection{Wind}

\subsection{Non-inertial frame corrections (Earth’s rotation)}



\newpage
\section{3DOF - Analytical analysis}
\subsection{Forces}
\subsection{Moment}
\subsection{Static stability}

\end{document}